\documentclass[../Odyssey_Project.tex]{subfiles}

\begin{document}
\setcounter{section}{15}
\section{Induction}

\subsubsection*{Overview}
\begin{itemize}
    \item Induced modules play a critical role in the representation theory of finite groups.
    \item Two equivalent definitions of induction are presented: one via tensor products over arbitrary rings (slick but heavier), the other using only tensor products over fields (clumsier but more elementary). A third definition is developed in Exercise~\ref{ex:16.1}.
    \item Let $H \leq G$ and let $V$ be an $FH$-module. Since the group algebra $FG$ is an $(FG, FH)$-bimodule, we form the $FG$-module $FG \otimes_{FH} V$, the \emph{induced} $FG$-module of $V$ (or the \emph{induction} of $V$ to $G$), denoted $\mathrm{Ind}_H^G V$. Other notations: $V^G$, $V \uparrow G$, and $V \uparrow_H^G$.
    \item For an $FG$-module $U$, regarding $U$ as an $FH$-module yields the \emph{restriction} $\mathrm{Res}_H^G U$. Other notations: $U_H$, $U \downarrow H$, and $U \downarrow_H^G$.
\end{itemize}

\begin{lemma}
\label{lem:16.1}
We have $\dim_F \mathrm{Ind}_H^G V = |G : H| \dim_F V$ for $V$ an $FH$-module. Moreover, as $F$-vector spaces we have
\[
\mathrm{Ind}_H^G V = \bigoplus_{t \in T} t \otimes V,
\]
where $T$ is a transversal for $H$ in $G$ and $t \otimes V = \{t \otimes v \mid v \in V\}$.
\end{lemma}

\begin{proof}
    We take $T$ to be a left transversal, so that
    \[
        G=\bigsqcup_{t\in T}tH.
    \]
    Consequently,
    \[
        FG=\bigoplus_{t\in T}tFH
    \]
    as right $FH$-modules. For each $t\in T$, the map $tFH\to FH$ sending $th$ to $h$ is an isomorphism of right $FH$-modules. Hence Propositions~\ref{prop:12.2} and~\ref{prop:12.3} give the following isomorphisms of $F$-vector spaces:
    \[
        \mathrm{Ind}_H^G V
        =
        FG\otimes_{FH}V
        \cong
        \bigoplus_{t\in T}\bigl(tFH\otimes_{FH}V\bigr)
        \cong
        \bigoplus_{t\in T}V.
    \]
    Under the last isomorphism, the copy of $V$ indexed by $t$ is precisely the subspace
    \[
        t\otimes V=\{t\otimes v\mid v\in V\}.
    \]
    Indeed, if $g=th$ with $t\in T$ and $h\in H$, then the balancing relation in the tensor product gives
    \[
        g\otimes v=th\otimes v=t\otimes hv\in t\otimes V.
    \]
    Since the preceding decomposition is a direct sum, we obtain
    \[
        \mathrm{Ind}_H^G V=\bigoplus_{t\in T}t\otimes V
    \]
    as $F$-vector spaces. Finally, $|T|=|G : H|$, so
    \[
        \dim_F\mathrm{Ind}_H^G V
        =
        |T|\dim_FV
        =
        |G : H|\dim_FV.
    \]
\end{proof}

\begin{namedtheorem}[Frobenius Reciprocity Theorem]
\label{thm:frobenius-reciprocity}
Let $U$ be an $FG$-module, let $H \leq G$, and let $V$ be an $FH$-module. Then as $F$-vector spaces we have $\mathrm{Hom}_{FH}(V, \mathrm{Res}_H^G U) \cong \mathrm{Hom}_{FG}(\mathrm{Ind}_H^G V, U)$.
\end{namedtheorem}

\begin{proof}
    Let $\varphi\in\Hom_{FH}(V,\mathrm{Res}_H^G U)$. Define
    \[
        f_\varphi\colon FG\times V\to U,
        \qquad
        f_\varphi(a,v)=a\varphi(v).
    \]
    This map is additive in each variable. Moreover, if $b\in FH$, then
    \[
        f_\varphi(ab,v)
        =
        ab\varphi(v)
        =
        a\varphi(bv)
        =
        f_\varphi(a,bv),
    \]
    because $\varphi$ is $FH$-linear. Also, for $c\in FG$,
    \[
        f_\varphi(ca,v)=ca\varphi(v)=c f_\varphi(a,v).
    \]
    Thus $f_\varphi$ is balanced, so the universal property of the tensor product gives a unique $FG$-module homomorphism
    \[
        \Gamma(\varphi)\colon FG\otimes_{FH}V\to U
    \]
    satisfying
    \[
        \Gamma(\varphi)(a\otimes v)=a\varphi(v).
    \]
    This defines an $F$-linear map
    \[
        \Gamma\colon
        \Hom_{FH}(V,\mathrm{Res}_H^G U)
        \longrightarrow
        \Hom_{FG}(\mathrm{Ind}_H^G V,U).
    \]
    Conversely, let $\theta\in\Hom_{FG}(\mathrm{Ind}_H^G V,U)$, and define
    \[
        \Delta(\theta)\colon V\to\mathrm{Res}_H^G U,
        \qquad
        \Delta(\theta)(v)=\theta(1\otimes v).
    \]
    For $h\in H$ and $v\in V$, the balancing relation gives
    \[
        \Delta(\theta)(hv)
        =
        \theta(1\otimes hv)
        =
        \theta(h\otimes v)
        =
        h\theta(1\otimes v)
        =
        h\Delta(\theta)(v).
    \]
    Hence $\Delta(\theta)$ is $FH$-linear. The two constructions are inverse: for $v\in V$,
    \[
        \Delta(\Gamma(\varphi))(v)
        =
        \Gamma(\varphi)(1\otimes v)
        =
        \varphi(v),
    \]
    while for every elementary tensor $a\otimes v$,
    \[
        \Gamma(\Delta(\theta))(a\otimes v)
        =
        a\theta(1\otimes v)
        =
        \theta(a\otimes v).
    \]
    Since elementary tensors span $\mathrm{Ind}_H^G V$, we have $\Gamma\Delta(\theta)=\theta$. Therefore $\Gamma$ is the required isomorphism of $F$-vector spaces.
\end{proof}

\subsubsection*{Induced Characters}
\begin{itemize}
    \item Restrict attention to $F = \C$. If $V$ is a $\C H$-module with character $\varphi$, the character of $\mathrm{Ind}_H^G V$ is denoted $\varphi^G$ and called an \emph{induced character}; we have $\varphi^G(1) = |G : H| \varphi(1)$.
    \item If $U$ is a $\C G$-module with character $\chi$, the character of $\mathrm{Res}_H^G U$ is denoted $\chi|_H$, and $\chi|_H(1) = \chi(1)$.
    \item The form of ``Frobenius reciprocity'' actually known to Frobenius (Theorem~\ref{thm:16.2}) follows from \nameref{thm:frobenius-reciprocity} and Theorem~\ref{thm:14.12}.
\end{itemize}

\begin{theorem}
\label{thm:16.2}
Let $H \leq G$, let $U$ be a $\C G$-module having character $\chi$, and let $V$ be a $\C H$-module having character $\varphi$. Then $(\varphi, \chi|_H)_H = (\varphi^G, \chi)_G$.
\end{theorem}

\begin{proof}
    The character of $\mathrm{Res}_H^G U$ is $\chi|_H$, while the character of $\mathrm{Ind}_H^G V$ is $\varphi^G$. Applying Theorem~\ref{thm:14.12} first to $H$ and then to $G$ gives
    \[
        (\varphi,\chi|_H)_H
        =
        \dim_{\C}\Hom_{\C H}(V,\mathrm{Res}_H^G U)
    \]
    and
    \[
        (\varphi^G,\chi)_G
        =
        \dim_{\C}\Hom_{\C G}(\mathrm{Ind}_H^G V,U).
    \]
    These two Hom-spaces are isomorphic by \nameref{thm:frobenius-reciprocity}, so their dimensions are equal. Therefore
    \[
        (\varphi,\chi|_H)_H=(\varphi^G,\chi)_G.
    \]
\end{proof}

\subsubsection*{Induction--Restriction Tables}
\begin{itemize}
    \item If $H \leq G$ and $\varphi_1, \dots, \varphi_r$ and $\chi_1, \dots, \chi_s$ are the irreducible characters of $H$ and $G$ respectively, the \emph{induction-restriction table} of $(G, H)$ is the $r \times s$ matrix whose $(i,j)$-entry gives the common value of $(\varphi_i^G, \chi_j)_G$ and $(\varphi_i, \chi_j|_H)_H$.
    \item By Corollary~\ref{cor:15.4}, the $i$th row gives the multiplicities with which the $\chi_j$ appear in the decomposition of $\varphi_i^G$, while the $j$th column gives the multiplicities with which the $\varphi_i$ appear in the decomposition of $\chi_j|_H$.
    \item One can compute the table by calculating the restrictions of the $\chi_j$ and reading off the expressions for the $\varphi_i^G$ in terms of the $\chi_j$.
    \item For example, with $G = S_3$ and $H = S_2 \leq G$, inspection gives $\chi_1|_H = \varphi_1$, $\chi_2|_H = \varphi_2$, $\chi_3|_H = \varphi_1 + \varphi_2$, so the induction-restriction table yields $\varphi_1^G = \chi_1 + \chi_3$ and $\varphi_2^G = \chi_2 + \chi_3$.
\end{itemize}

\begin{proposition}
\label{prop:16.3}
Let $H \leq G$, and let $V$ be a $\C H$-module having character $\chi$. If $T$ is a transversal for $H$ in $G$, then for any $g \in G$ we have
\[
\chi^G(g) = \sum_{\substack{t \in T, \\ t^{-1} g t \in H}} \chi(t^{-1} g t).
\]
\end{proposition}

\begin{proof}
    Choose a basis $\{v_1,\ldots,v_d\}$ of $V$. By Lemma~\ref{lem:16.1},
    \[
        \{t\otimes v_i\mid t\in T,\ 1\leq i\leq d\}
    \]
    is a basis of $\mathrm{Ind}_H^G V$. Fix $g\in G$. For each $t\in T$, there are unique $s\in T$ and $h\in H$ such that $gt=sh$. Therefore
    \[
        g(t\otimes v)=gt\otimes v=sh\otimes v=s\otimes hv,
    \]
    so $g$ maps the subspace $t\otimes V$ into $s\otimes V$.\\
    If $s\neq t$, then this map contributes nothing to the trace, because it sends the whole $t\otimes V$ block into a different block. The equality $s=t$ holds precisely when
    \[
        gt=th
        \quad\Longleftrightarrow\quad
        h=t^{-1}gt\in H.
    \]
    In this case,
    \[
        g(t\otimes v)=t\otimes(t^{-1}gt)v.
    \]
    Under the isomorphism $V\to t\otimes V$ sending $v$ to $t\otimes v$, the action of $g$ on this block is therefore the same as the action of $t^{-1}gt$ on $V$. Its contribution to the trace is $\chi(t^{-1}gt)$. Summing the contributions of all diagonal blocks gives
    \[
        \chi^G(g)
        =
        \sum_{\substack{t\in T,\\t^{-1}gt\in H}}
        \chi(t^{-1}gt).
    \]
\end{proof}

\begin{corollary}
\label{cor:16.4}
Let $\chi$ be a character of $H \leq G$, and let $g \in G$. Then
\[
\chi^G(g) = \frac{1}{|H|} \sum_{\substack{x \in G, \\ x^{-1} g x \in H}} \chi(x^{-1} g x).
\]
\end{corollary}

\begin{proof}
    Let $T$ be a left transversal for $H$ in $G$. Every $x\in G$ can be written uniquely as $x=th$ with $t\in T$ and $h\in H$. For such an $x$,
    \[
        x^{-1}gx=h^{-1}(t^{-1}gt)h.
    \]
    Hence $x^{-1}gx\in H$ if and only if $t^{-1}gt\in H$. Whenever this occurs, $\chi$ being a class function on $H$ gives
    \[
        \chi(x^{-1}gx)
        =
        \chi\bigl(h^{-1}(t^{-1}gt)h\bigr)
        =
        \chi(t^{-1}gt).
    \]
    Thus, for each $t\in T$ with $t^{-1}gt\in H$, the $|H|$ elements $th$, where $h\in H$, all make the same contribution. Therefore
    \[
        \sum_{\substack{x\in G,\\x^{-1}gx\in H}}\chi(x^{-1}gx)
        =
        |H|
        \sum_{\substack{t\in T,\\t^{-1}gt\in H}}\chi(t^{-1}gt).
    \]
    Proposition~\ref{prop:16.3} now gives
    \[
        \chi^G(g)
        =
        \frac{1}{|H|}
        \sum_{\substack{x\in G,\\x^{-1}gx\in H}}\chi(x^{-1}gx).
    \]
\end{proof}

\begin{proposition}
\label{prop:16.5}
Let $\chi$ be a character of $H \leq G$, and let $g \in G$, and let $s$ be the number of conjugacy classes of $H$ whose members are conjugate in $G$ to $g$. If $s = 0$, then $\chi^G(g) = 0$. Otherwise, if we let $h_1, \dots, h_s$ be representatives of these $s$ conjugacy classes of $H$, then
\[
\chi^G(g) = \sum_{i=1}^s \frac{|C_G(g)|}{|C_H(h_i)|} \chi(h_i).
\]
\end{proposition}

\begin{proof}
    If $s=0$, no conjugate of $g$ lies in $H$. Hence the sum in Corollary~\ref{cor:16.4} is empty, and therefore $\chi^G(g)=0$.\\
    Now assume $s>0$. For each $1\leq i\leq s$, let
    \[
        X_i
        =
        \{x\in G\mid x^{-1}gx\text{ is conjugate in }H\text{ to }h_i\}.
    \]
    The sets $X_i$ are pairwise disjoint, and their union is $\{x\in G\mid x^{-1}gx\in H\}$. Since $\chi$ is constant on each conjugacy class of $H$, Corollary~\ref{cor:16.4} gives
    \[
        \chi^G(g)
        =
        \frac{1}{|H|}
        \sum_{i=1}^{s}|X_i|\chi(h_i).
    \]
    It remains to calculate $|X_i|$. Choose $t_i\in G$ such that $t_i^{-1}gt_i=h_i$. We claim that
    \[
        X_i=C_G(g)t_iH.
    \]
    Indeed, if $c\in C_G(g)$ and $h\in H$, then
    \[
        (ct_ih)^{-1}g(ct_ih)
        =
        h^{-1}t_i^{-1}c^{-1}gct_ih
        =
        h^{-1}h_ih,
    \]
    so $ct_ih\in X_i$. This proves $C_G(g)t_iH\subseteq X_i$. Conversely, if $x\in X_i$, then for some $h\in H$,
    \[
        x^{-1}gx
        =
        h^{-1}h_ih
        =
        h^{-1}t_i^{-1}gt_ih.
    \]
    Put $c=xh^{-1}t_i^{-1}$. Then
    \[
        c^{-1}gc
        =
        t_ihx^{-1}gxh^{-1}t_i^{-1}
        =
        g,
    \]
    so $c\in C_G(g)$ and $x=ct_ih\in C_G(g)t_iH$. Thus the claim follows.\\
    Proposition~\ref{prop:3.15} now gives
    \[
        |X_i|
        =
        \frac{|C_G(g)|\,|H|}
        {|t_i^{-1}C_G(g)t_i\cap H|}.
    \]
    By Exercise~\ref{ex:3.7},
    \[
        t_i^{-1}C_G(g)t_i
        =
        C_G(t_i^{-1}gt_i)
        =
        C_G(h_i),
    \]
    and hence $t_i^{-1}C_G(g)t_i\cap H=C_H(h_i)$. Therefore
    \[
        \frac{|X_i|}{|H|}
        =
        \frac{|C_G(g)|}{|C_H(h_i)|}.
    \]
    Substituting this into the earlier expression for $\chi^G(g)$ yields
    \[
        \chi^G(g)
        =
        \sum_{i=1}^{s}
        \frac{|C_G(g)|}{|C_H(h_i)|}\chi(h_i).
    \]
\end{proof}

\begin{corollary}
\label{cor:16.6}
Let $\chi$ be a character of $H \leq G$. Let $g \in G$, and suppose that the number $s$ of conjugacy classes of $H$ whose members are conjugate in $G$ to $g$ is positive. Let $\ell$ be the order of the conjugacy class in $G$ of $g$; let $h_1, \dots, h_s$ be representatives of these $s$ conjugacy classes of $H$, and let $k_1, \dots, k_s$ be the orders of these classes. Then
\[
\chi^G(g) = \sum_{i=1}^s |G : H| \frac{k_i}{\ell} \chi(h_i).
\]
\end{corollary}

\begin{proof}
    Proposition~\ref{prop:3.13} gives
    \[
        \ell=|G : C_G(g)|
        \qquad\text{and}\qquad
        k_i=|H : C_H(h_i)|.
    \]
    Hence
    \[
        |C_G(g)|=\frac{|G|}{\ell},
        \qquad
        |C_H(h_i)|=\frac{|H|}{k_i}.
    \]
    It follows that
    \[
        \frac{|C_G(g)|}{|C_H(h_i)|}
        =
        \frac{|G|/\ell}{|H|/k_i}
        =
        |G : H|\frac{k_i}{\ell}.
    \]
    Substituting this equality into Proposition~\ref{prop:16.5} gives
    \[
        \chi^G(g)
        =
        \sum_{i=1}^{s}|G : H|\frac{k_i}{\ell}\chi(h_i).
    \]
\end{proof}

\subsubsection*{Frobenius Groups}
\begin{itemize}
    \item We end the section by presenting an important theorem on finite groups which, quite remarkably, has no known proof that does not use character theory.
    \item A \emph{Frobenius group} is a group $G$ having a subgroup $H$ such that $H \cap g H g^{-1} = 1$ for every $g \in G - H$. We call $H$ a \emph{Frobenius complement} of $G$.
    \item If $G$ is a finite Frobenius group, applying \nameref{thm:frobenius} yields a normal subgroup $N$, called the \emph{Frobenius kernel} of $G$. One has $G = N \rtimes H$, $N \cap H = 1$, and $|N||H| = |G|$.
    \item For $G = N \rtimes H$ a finite Frobenius group with Frobenius kernel $N$ and complement $H$, the conjugation map $\varphi \colon H \to \Aut(N)$ has the property that for each $1 \neq h \in H$, $\varphi(h)$ is a \emph{fixed-point-free automorphism} of $N$: $\varphi(h)(n) = n$ forces $n = 1$. In particular, $\varphi$ is injective.
    \item Conversely, given finite groups $N$ and $H$ and a homomorphism $\varphi \colon H \to \Aut(N)$ for which every $\varphi(h)$, $h \neq 1$, is fixed-point-free, $N \rtimes_\varphi H$ is a Frobenius group.
    \item Example: if $p, q$ are distinct primes with $p \equiv 1 \pmod{q}$, then any monomorphism $\varphi \colon \Z_q \to \Aut(\Z_p)$ is fixed-point-free, and hence the unique non-abelian group of order $pq$ is a Frobenius group.
    \item A theorem of Thompson (his 1959 Chicago Ph.D. thesis) states that any finite group having a fixed-point-free automorphism of prime order is nilpotent. Hence the Frobenius kernel of a finite Frobenius group is nilpotent.
\end{itemize}

\begin{namedtheorem}[Frobenius' Theorem]
\label{thm:frobenius}
Let $G$ be a transitive permutation group on a finite set $X$, and suppose that each non-identity element of $G$ fixes at most one element of $X$. Then the union of the identity element and the set of elements of $G$ that have no fixed points is a normal subgroup of $G$.
\end{namedtheorem}

\begin{proof}
    Let
    \[
        N=\{1\}\cup\{g\in G\mid gx\neq x\text{ for every }x\in X\}.
    \]
    Thus $N$ consists of the identity together with all fixed-point-free elements. We must prove that $N\trianglelefteq G$. Put $n=|X|$, fix $x_0\in X$, and let $H=G_{x_0}$. Since $X$ is transitive, Proposition~\ref{prop:3.4} gives $X\cong G/H$, and hence
    \[
        n=|G : H|,
        \qquad
        |G|=n|H|.
    \]
    Also, Lemma~\ref{lem:3.2} shows that every point stabilizer is conjugate to $H$, so every point stabilizer has order $|H|$.\\
    Consider the set
    \[
        \mathcal{P}
        =
        \{(g,x)\in(G-\{1\})\times X\mid gx=x\}.
    \]
    For each $x\in X$, exactly $|G_x|-1=|H|-1$ non-identity elements fix $x$. Therefore
    \[
        |\mathcal{P}|=n(|H|-1).
    \]
    By hypothesis, each non-identity element of $G$ fixes at most one point. Thus every non-identity element which has a fixed point occurs in exactly one pair in $\mathcal{P}$. These are precisely the elements of $G-N$, so
    \[
        |G-N|=n(|H|-1).
    \]
    Consequently,
    \[
        |N|
        =
        |G|-n(|H|-1)
        =
        n|H|-n(|H|-1)
        =
        n.
    \]
    Notice also that conjugation preserves the number of fixed points: $aga^{-1}$ fixes $x$ if and only if $g$ fixes $a^{-1}x$. Hence $N$ is a union of conjugacy classes, although we do not yet know that it is a subgroup.\\
    We next determine how conjugacy in $G$ behaves on $H-\{1\}$. Let $1\neq h\in H$, and suppose
    \[
        h=ah'a^{-1}
    \]
    for some $a\in G$ and $h'\in H$. Since $h\in H$, it fixes $x_0$. Since $h'$ fixes $x_0$, its conjugate $ah'a^{-1}=h$ fixes $ax_0$. Thus the non-identity element $h$ fixes both $x_0$ and $ax_0$. The hypothesis forces $ax_0=x_0$, so $a\in H$. Therefore two non-identity elements of $H$ are conjugate in $G$ only if they are already conjugate in $H$.\\
    Taking $h'=h$ in the same argument shows that every element of $C_G(h)$ lies in $H$. Hence
    \[
        C_G(h)=C_H(h)
    \]
    for every $1\neq h\in H$. Finally, if $g\notin N$, then $g$ fixes a point $x=ax_0$ for some $a\in G$. Thus
    \[
        a^{-1}ga\in G_{x_0}=H.
    \]
    Since $g\neq1$, every element outside $N$ is therefore conjugate to a non-identity element of $H$.\\
    Let $\varphi_1,\ldots,\varphi_s$ be the irreducible characters of $H$, with $\varphi_1$ the principal character, and let $\chi_1$ be the principal character of $G$. For each $i$, consider the induced character $\varphi_i^G$. Lemma~\ref{lem:16.1} gives
    \[
        \varphi_i^G(1)=|G : H|\varphi_i(1)=n\varphi_i(1).
    \]
    If $1\neq h\in H$, then the preceding conjugacy argument and Proposition~\ref{prop:16.5} give
    \[
        \varphi_i^G(h)
        =
        \frac{|C_G(h)|}{|C_H(h)|}\varphi_i(h)
        =
        \varphi_i(h).
    \]
    If $y\in N-\{1\}$, then no conjugate of $y$ lies in $H$, because every non-identity element of $H$ has a fixed point. Hence Proposition~\ref{prop:16.5} gives $\varphi_i^G(y)=0$.\\
    For each $2\leq i\leq s$, define the virtual character
    \[
        \chi_i
        =
        \varphi_i^G-\varphi_i(1)\varphi_1^G+\varphi_i(1)\chi_1.
    \]
    Using the values just calculated, we obtain
    \[
        \begin{array}{c|ccc}
            &1&1\neq h\in H&y\in N-\{1\}\\
            \hline
            \chi_i&\varphi_i(1)&\varphi_i(h)&\varphi_i(1)
        \end{array}.
    \]
    We now show that $\chi_i$ is an irreducible character. The hypothesis gives the disjoint partition
    \[
        G
        =
        N\sqcup\bigsqcup_{x\in X}(G_x-\{1\}),
    \]
    because every non-identity element outside $N$ fixes exactly one point. Therefore
    \begin{align*}
        (\chi_i,\chi_i)_G
        &=
        \frac{1}{|G|}
        \left(
            \sum_{g\in N}|\chi_i(g)|^2
            +
            \sum_{x\in X}\sum_{1\neq g\in G_x}|\chi_i(g)|^2
        \right) \\
        &=
        \frac{1}{|G|}
        \left(
            n\varphi_i(1)^2
            +
            n\sum_{1\neq h\in H}|\varphi_i(h)|^2
        \right) \\
        &=
        \frac{1}{|H|}
        \sum_{h\in H}|\varphi_i(h)|^2 \\
        &=
        (\varphi_i,\varphi_i)_H
        =
        1,
    \end{align*}
    where the last equality follows from \nameref{thm:row-orthogonality}, applied to $H$. Since $\chi_i$ is a virtual character, Corollary~\ref{cor:15.2} says that its inner product with itself is the sum of the squares of its integer coefficients in the irreducible-character basis. Since this sum is $1$, either $\chi_i$ or $-\chi_i$ is irreducible. But
    \[
        \chi_i(1)=\varphi_i(1)>0,
    \]
    so $\chi_i$, rather than $-\chi_i$, is an irreducible character of $G$.\\
    Finally, define the character
    \[
        \chi=\sum_{i=1}^{s}\varphi_i(1)\chi_i,
    \]
    where $\chi_1$ is the principal character and the characters $\chi_i$ for $i\geq2$ are those constructed above. If $y\in N$, then $\chi_i(y)=\varphi_i(1)$ for every $i$; for $i=1$, both sides are $1$. Thus Corollary~\ref{cor:14.2}, applied to $H$, gives
    \[
        \chi(y)
        =
        \sum_{i=1}^{s}\varphi_i(1)^2
        =
        |H|.
    \]
    If $g\notin N$, then $g$ is conjugate in $G$ to some $1\neq h\in H$. Since characters are class functions,
    \[
        \chi(g)
        =
        \chi(h)
        =
        \sum_{i=1}^{s}\varphi_i(1)\varphi_i(h).
    \]
    Apply \nameref{thm:col-orthogonality} to the columns of the character table of $H$ represented by $1$ and $h^{-1}$. Using part (iv) of Proposition~\ref{prop:14.4}, we obtain
    \[
        \sum_{i=1}^{s}\varphi_i(1)\varphi_i(h)
        =
        \sum_{i=1}^{s}\varphi_i(1)\overline{\varphi_i(h^{-1})}
        =
        0.
    \]
    Hence
    \[
        \chi(g)
        =
        \begin{cases}
            |H|,&g\in N,\\
            0,&g\notin N.
        \end{cases}
    \]
    In particular, $\chi(1)=|H|$, and therefore
    \[
        N=\{g\in G\mid\chi(g)=\chi(1)\}.
    \]
    Part (vi) of Proposition~\ref{prop:14.4} says that this set is a normal subgroup of $G$. Thus $N\trianglelefteq G$.
\end{proof}

\subsection*{Exercises}

\noindent Throughout these exercises, $G$ is a finite group, $H$ is a subgroup of $G$, $F$ is a field, and all $FG$-modules are finitely generated.

\begin{exercise}
\label{ex:16.1}
Let $V$ be an $FH$-module, and let $W$ be an $FG$-module containing $V$. Suppose that $W$ has the property that for any $FG$-module $U$ and any $\varphi \in \mathrm{Hom}_{FH}(V, U)$, there is a unique $FG$-module homomorphism from $W$ to $U$ whose restriction to $V$ is $\varphi$. (In this case, we say that $W$ is relatively $H$-free with respect to $V$; compare with Exercises~\ref{ex:13.1}--\ref{ex:13.2}.) Show that $W \cong \mathrm{Ind}_H^G V$.
\end{exercise}
\begin{solution}
    Let
    \[
        I=\mathrm{Ind}_H^G V=FG\otimes_{FH}V.
    \]
    By Lemma~\ref{lem:16.1}, the map
    \[
        \eta\colon V\to I,
        \qquad
        \eta(v)=1\otimes v,
    \]
    identifies $V$ with the $FH$-submodule $1\otimes V$ of $I$. Apply the assumed property of $W$ with $U=I$ and $\varphi=\eta$. We obtain an $FG$-module homomorphism
    \[
        \alpha\colon W\to I
    \]
    such that $\alpha(v)=1\otimes v$ for every $v\in V$.\\
    In the other direction, define
    \[
        \beta\colon I\to W,
        \qquad
        \beta(g\otimes v)=gv.
    \]
    This is well-defined because, for $h\in H$,
    \[
        \beta(gh\otimes v)=ghv=g(hv)=\beta(g\otimes hv),
    \]
    and it is $FG$-linear. For $v\in V$, we have
    \[
        (\beta\alpha)(v)=\beta(1\otimes v)=v.
    \]
    Thus both $\beta\alpha$ and $\operatorname{id}_W$ are $FG$-module homomorphisms from $W$ to $W$ whose restriction to $V$ is the identity. The uniqueness in the defining property of $W$ therefore gives
    \[
        \beta\alpha=\operatorname{id}_W.
    \]
    Also,
    \[
        (\alpha\beta)(1\otimes v)=\alpha(v)=1\otimes v.
    \]
    Since $I$ is generated as an $FG$-module by $1\otimes V$, this implies that $\alpha\beta=\operatorname{id}_I$. Hence $\alpha$ and $\beta$ are inverse $FG$-module isomorphisms, and so
    \[
        W\cong\mathrm{Ind}_H^G V.
    \]
\end{solution}

\begin{exercise}
\label{ex:16.2}
Let $W$ be an $FG$-module that is generated by an $FH$-submodule $V$, and suppose that $\dim_F W = |G : H| \dim_F V$. Show that we must have $W \cong \mathrm{Ind}_H^G V$.
\end{exercise}
\begin{solution}
    Define
    \[
        \pi\colon FG\otimes_{FH}V\to W,
        \qquad
        \pi(g\otimes v)=gv.
    \]
    This map is well-defined because
    \[
        \pi(gh\otimes v)=ghv=g(hv)=\pi(g\otimes hv)
    \]
    for every $g\in G$, $h\in H$, and $v\in V$, and it is $FG$-linear. Since $W$ is generated as an $FG$-module by $V$, every element of $W$ is a sum of elements of the form $gv$. Hence $\pi$ is surjective.\\
    By Lemma~\ref{lem:16.1},
    \[
        \dim_F\mathrm{Ind}_H^G V
        =
        |G : H|\dim_FV
        =
        \dim_FW.
    \]
    Thus $\pi$ is a surjective linear map between finite-dimensional vector spaces of the same dimension, so it is also injective. Therefore
    \[
        W\cong\mathrm{Ind}_H^G V
    \]
    as $FG$-modules.
\end{solution}

\begin{exercise}
\label{ex:16.3}
Show that $\mathrm{Hom}_{FG}(U, \mathrm{Ind}_H^G V) \cong \mathrm{Hom}_{FH}(\mathrm{Res}_H^G U, V)$ as $F$-vector spaces for any $FG$-module $U$ and $FH$-module $V$. (\textsc{Hint}: Show first that if $\varphi \in \mathrm{Hom}_{FH}(\mathrm{Res}_H^G U, V)$, then the map sending $u \in U$ to $\sum_{t \in T} t \otimes \varphi(t^{-1} u)$, where $T$ is a transversal for $H$ in $G$, lies in $\mathrm{Hom}_{FG}(U, \mathrm{Ind}_H^G V)$.)
\end{exercise}
\begin{solution}
    Choose a transversal $T$ for $H$ in $G$ which contains $1$. For
    \[
        \varphi\in\Hom_{FH}(\mathrm{Res}_H^G U,V),
    \]
    define
    \[
        \Phi(\varphi)(u)
        =
        \sum_{t\in T}t\otimes\varphi(t^{-1}u).
    \]
    We first check that $\Phi(\varphi)$ is $FG$-linear. Let $x\in G$. For each $t\in T$, write $xt=sh$ with $s\in T$ and $h\in H$. As $t$ runs through $T$, so does $s$. Since $\varphi$ is $FH$-linear and $ht^{-1}=s^{-1}x$, we obtain
    \begin{align*}
        x\Phi(\varphi)(u)
        &=
        \sum_{t\in T}xt\otimes\varphi(t^{-1}u) \\
        &=
        \sum_{t\in T}s\otimes h\varphi(t^{-1}u) \\
        &=
        \sum_{t\in T}s\otimes\varphi(ht^{-1}u) \\
        &=
        \sum_{s\in T}s\otimes\varphi(s^{-1}xu)
        =
        \Phi(\varphi)(xu).
    \end{align*}
    Hence
    \[
        \Phi\colon
        \Hom_{FH}(\mathrm{Res}_H^G U,V)
        \longrightarrow
        \Hom_{FG}(U,\mathrm{Ind}_H^G V)
    \]
    is a linear map.\\
    By Lemma~\ref{lem:16.1}, every $z\in\mathrm{Ind}_H^G V$ has a unique expression
    \[
        z=\sum_{t\in T}t\otimes v_t.
    \]
    Let
    \[
        p\colon\mathrm{Ind}_H^G V\to V,
        \qquad
        p(z)=v_1
    \]
    be the projection onto the $1\otimes V$ component. This map is $FH$-linear: multiplication by an element of $H$ preserves $1\otimes V$, and it cannot move any $t\otimes V$ with $t\neq1$ into $1\otimes V$. Thus, for
    \[
        f\in\Hom_{FG}(U,\mathrm{Ind}_H^G V),
    \]
    the map $\Psi(f)=p\circ f$ lies in $\Hom_{FH}(\mathrm{Res}_H^G U,V)$.\\
    For $\varphi\in\Hom_{FH}(\mathrm{Res}_H^G U,V)$, the $1\otimes V$ component of $\Phi(\varphi)(u)$ is $1\otimes\varphi(u)$. Therefore $\Psi\Phi(\varphi)=\varphi$. Conversely, suppose that
    \[
        f(u)=\sum_{t\in T}t\otimes v_t.
    \]
    Since $f$ is $FG$-linear, $f(t^{-1}u)=t^{-1}f(u)$. The $1\otimes V$ component of $t^{-1}f(u)$ is $v_t$, because $t^{-1}s\in H$ for $s,t\in T$ exactly when $s=t$. Hence
    \[
        \Phi\Psi(f)(u)
        =
        \sum_{t\in T}t\otimes v_t
        =
        f(u).
    \]
    Thus $\Phi$ and $\Psi$ are inverse linear maps, and therefore
    \[
        \Hom_{FG}(U,\mathrm{Ind}_H^G V)
        \cong
        \Hom_{FH}(\mathrm{Res}_H^G U,V)
    \]
    as $F$-vector spaces.
\end{solution}

\begin{exercise}
\label{ex:16.4}
Show that if $X$ is a transitive $G$-set, then $FX \cong \mathrm{Ind}_{G_x}^G F$ for any $x \in X$.
\end{exercise}
\begin{solution}
    Fix $x\in X$, put $H=G_x$, and regard $F$ as the trivial $FH$-module. Define
    \[
        \Phi\colon FG\otimes_{FH}F\to FX,
        \qquad
        \Phi(g\otimes\lambda)=\lambda(gx).
    \]
    This is well-defined because if $h\in H$, then $hx=x$, and hence
    \[
        \Phi(gh\otimes\lambda)
        =
        \lambda(ghx)
        =
        \lambda(gx)
        =
        \Phi(g\otimes h\lambda).
    \]
    The map is also $FG$-linear. Since $X$ is transitive, every element of $X$ has the form $gx$, so $\Phi$ is surjective.\\
    By the orbit-stabilizer theorem and Lemma~\ref{lem:16.1},
    \[
        \dim_FFX
        =
        |X|
        =
        |G : G_x|
        =
        \dim_F\mathrm{Ind}_{G_x}^G F.
    \]
    Thus $\Phi$ is an isomorphism, and therefore
    \[
        FX\cong\mathrm{Ind}_{G_x}^G F.
    \]
\end{solution}

\begin{exercise}
\label{ex:16.5}
Compute the induction-restriction table of $(A_5, A_4)$.
\end{exercise}
\begin{solution}
    Take $H=A_4$ to be the subgroup of $A_5$ fixing $5$. Let
    \[
        d=(1\,2)(3\,4),
        \qquad
        c=(1\,2\,3),
    \]
    and let $\omega$ be a primitive cube root of unity. The conjugacy classes of $H$ have representatives $1,d,c,c^{-1}$ and respective orders $1,3,4,4$. Its irreducible character table is
    \[
        \begin{array}{c|rrrr}
            &1&d&c&c^{-1}\\
            \hline
            \varphi_1&1&1&1&1\\
            \varphi_2&1&1&\omega&\omega^2\\
            \varphi_3&1&1&\omega^2&\omega\\
            \varphi_4&3&-1&0&0
        \end{array}.
    \]
    We use the numbering of the irreducible characters of $A_5$ from Example~\ref{ex:15.eg.9}, so their degrees are
    \[
        \chi_1(1)=1,
        \quad
        \chi_2(1)=4,
        \quad
        \chi_3(1)=5,
        \quad
        \chi_4(1)=\chi_5(1)=3.
    \]
    Restricting the $A_5$-table to elements of $H$ gives
    \[
        \begin{array}{c|rrrr}
            &1&d&c&c^{-1}\\
            \hline
            \chi_1|_H&1&1&1&1\\
            \chi_2|_H&4&0&1&1\\
            \chi_3|_H&5&1&-1&-1\\
            \chi_4|_H&3&-1&0&0\\
            \chi_5|_H&3&-1&0&0
        \end{array}.
    \]
    Comparing these rows with the $A_4$-table, and using $\omega+\omega^2=-1$, we obtain
    \[
        \begin{aligned}
            \chi_1|_H&=\varphi_1,\\
            \chi_2|_H&=\varphi_1+\varphi_4,\\
            \chi_3|_H&=\varphi_2+\varphi_3+\varphi_4,\\
            \chi_4|_H&=\varphi_4,\\
            \chi_5|_H&=\varphi_4.
        \end{aligned}
    \]
    Therefore the induction-restriction table of $(A_5,A_4)$, with rows indexed by the $\varphi_i$ and columns indexed by the $\chi_j$, is
    \[
        \begin{array}{c|ccccc}
            &\chi_1&\chi_2&\chi_3&\chi_4&\chi_5\\
            \text{degree}&1&4&5&3&3\\
            \hline
            \varphi_1&1&1&0&0&0\\
            \varphi_2&0&0&1&0&0\\
            \varphi_3&0&0&1&0&0\\
            \varphi_4&0&1&1&1&1
        \end{array}.
    \]
    Equivalently, Frobenius reciprocity gives
    \[
        \begin{aligned}
            \varphi_1^{A_5}&=\chi_1+\chi_2,\\
            \varphi_2^{A_5}&=\chi_3,\\
            \varphi_3^{A_5}&=\chi_3,\\
            \varphi_4^{A_5}&=\chi_2+\chi_3+\chi_4+\chi_5.
        \end{aligned}
    \]
    Since $|A_5 : A_4|=5$, the induced degrees must be $5,5,5,15$. The displayed decompositions have degrees
    \[
        1+4,
        \qquad
        5,
        \qquad
        5,
        \qquad
        4+5+3+3,
    \]
    respectively, which checks every row of the table.
\end{solution}

\begin{exercise}
\label{ex:16.6}
Compute the character table of the unique non-abelian group of order $pq$, where $p$ and $q$ are distinct primes and $p \equiv 1 \pmod{q}$.
\end{exercise}
\begin{solution}
    Since $p\equiv1\pmod q$, we have $q\mid p-1$. Choose $r\in(\Z_p)^\times$ of order $q$. The group has a presentation
    \[
        G=\langle a,b\mid a^p=b^q=1,\ bab^{-1}=a^r\rangle
        =
        \langle a\rangle\rtimes\langle b\rangle.
    \]
    Put
    \[
        N=\langle a\rangle,
        \qquad
        m=\frac{p-1}{q}.
    \]
    Conjugation by $b$ sends $a^u$ to $a^{ru}$. Hence the non-zero elements of $\Z_p$ split into $m$ orbits of order $q$ under multiplication by $r$. Choose representatives $u_1,\ldots,u_m$ for these orbits. The conjugacy classes contained in $N-\{1\}$ are
    \[
        C_k=\{a^{u_kr^t}\mid0\leq t<q\},
        \qquad
        1\leq k\leq m,
    \]
    and each has order $q$.\\
    For $1\leq s<q$, conjugation by $a^w$ gives
    \[
        a^w(a^vb^s)a^{-w}
        =
        a^{v+(1-r^s)w}b^s.
    \]
    Since $r^s\not\equiv1\pmod p$, multiplication by $1-r^s$ is a bijection of $\Z_p$. Thus
    \[
        D_s=\{a^vb^s\mid0\leq v<p\}
    \]
    is one conjugacy class of order $p$. Distinct values of $s$ give distinct classes because their images in the abelian quotient $G/N$ are distinct. Hence $G$ has
    \[
        1+m+(q-1)=q+\frac{p-1}{q}
    \]
    conjugacy classes.\\
    Let
    \[
        \epsilon=e^{2\pi i/p},
        \qquad
        \zeta=e^{2\pi i/q}.
    \]
    Since $bab^{-1}a^{-1}=a^{r-1}$ generates $N$, we have $G'=N$. Thus the $q$ linear characters are the lifts of the irreducible characters of $G/N\cong\Z_q$. They are
    \[
        \lambda_j(a)=1,
        \qquad
        \lambda_j(b)=\zeta^j,
        \qquad
        0\leq j<q.
    \]
    For $1\leq i\leq m$, define a non-principal linear character $\theta_i$ of $N$ by $\theta_i(a)=\epsilon^{u_i}$, and put $\psi_i=\theta_i^G$. Proposition~\ref{prop:16.3} gives
    \[
        \psi_i(a^{u_k})
        =
        \sum_{t=0}^{q-1}\epsilon^{u_i u_k r^t},
        \qquad
        \psi_i(b^s)=0.
    \]
    Moreover,
    \[
        \psi_i|_N
        =
        \sum_{t=0}^{q-1}\theta_i^{\,b^t}.
    \]
    Here $\theta_i^{\,b^t}(n)=\theta_i(b^{-t}nb^t)$ for $n\in N$.
    These $q$ conjugate characters of $N$ are distinct because $r$ has order $q$. Therefore, by Theorem~\ref{thm:16.2},
    \[
        (\psi_i,\psi_i)_G
        =
        (\theta_i,\psi_i|_N)_N
        =
        1.
    \]
    Hence $\psi_i$ is irreducible by Corollary~\ref{cor:15.3}. Characters obtained from different orbits have disjoint sets of irreducible constituents upon restriction to $N$, so the $\psi_i$ are distinct.\\
    The complete character table is therefore
    \[
        \begin{array}{c|c|c|c}
            \text{class representative}
            &1
            &a^{u_k}\ (1\leq k\leq m)
            &b^s\ (1\leq s<q)\\
            \text{class order}
            &1&q&p\\
            \hline
            \lambda_j\ (0\leq j<q)
            &1&1&\zeta^{js}\\
            \psi_i\ (1\leq i\leq m)
            &q&
            \displaystyle\sum_{t=0}^{q-1}\epsilon^{u_i u_k r^t}
            &0
        \end{array}.
    \]
    Finally, the squares of the degrees of these characters sum to
    \[
        q\cdot1^2+m q^2
        =
        q+\frac{p-1}{q}q^2
        =
        pq
        =
        |G|.
    \]
    By Corollary~\ref{cor:14.2}, no further irreducible characters exist.
\end{solution}

\begin{exercise}
\label{ex:16.7}
Suppose that $G$ is a Frobenius group having Frobenius kernel $N$. Show that the irreducible characters of $G$ either are lifts of irreducible characters of $G/N$ or are induced from the non-principal irreducible characters of $N$. When is it true that the inductions to $G$ of two distinct non-principal irreducible characters of $N$ are equal?
\end{exercise}
\begin{solution}
    Choose a Frobenius complement $H$, so that
    \[
        G=N\rtimes H.
    \]
    We first show that $H$ acts freely on the set of non-principal irreducible characters of $N$. Consider the Frobenius action of $G$ on the cosets of $H$. If $1\neq n\in N$, then
    \[
        C_G(n)\leq N.
    \]
    Indeed, suppose that $x\in C_G(n)$ but $x\notin N$. By the definition of $N$, the element $x$ is not fixed-point-free; by the Frobenius property, it therefore has exactly one fixed coset, say $y$. Since $n$ has no fixed points, $ny\neq y$, while
    \[
        x(ny)=n(xy)=ny.
    \]
    Thus $x$ fixes both $y$ and $ny$, contradicting the Frobenius property.\\
    Now fix $1\neq h\in H$. Suppose that conjugation by $h$ fixed the $N$-conjugacy class of some $1\neq n\in N$. Then
    \[
        hnh^{-1}=xnx^{-1}
    \]
    for some $x\in N$, so $x^{-1}h\in C_G(n)$. But $x^{-1}h\notin N$, contradicting $C_G(n)\leq N$. Thus conjugation by $h$ fixes only the identity conjugacy class of $N$.\\
    Write $\operatorname{Irr}(N)$ for the set of irreducible characters of $N$, and let $\mathcal{X}_N$ be the character table of $N$. Conjugation by $h$ permutes both $\operatorname{Irr}(N)$ and the conjugacy classes of $N$. The same character-table argument as in Exercises~\ref{ex:15.4}--\ref{ex:15.5} gives permutation matrices $P_h$ and $Q_h$ satisfying
    \[
        P_h\mathcal{X}_N=\mathcal{X}_NQ_h.
    \]
    Since $\mathcal{X}_N$ is invertible, $P_h$ and $Q_h$ are similar and have the same trace. These traces count the irreducible characters and conjugacy classes fixed by $h$, respectively. Only the identity class is fixed, so exactly one irreducible character is fixed. The principal character $1_N$ is always fixed; hence no non-principal irreducible character of $N$ is fixed by $h$. Therefore $H$ acts freely on
    \[
        \operatorname{Irr}(N)-\{1_N\}.
    \]
    For $\theta\in\operatorname{Irr}(N)$ and $h\in H$, write
    \[
        \theta^h(n)=\theta(h^{-1}nh).
    \]
    If $\theta\neq1_N$, Proposition~\ref{prop:16.3}, applied with $N\trianglelefteq G$ and transversal $H$, gives
    \[
        \theta^G|_N=\sum_{h\in H}\theta^h.
    \]
    Since the action is free, the characters in this sum are distinct and $\theta$ occurs exactly once. Frobenius reciprocity now gives
    \[
        (\theta^G,\theta^G)_G
        =
        (\theta,\theta^G|_N)_N
        =
        1.
    \]
    Hence $\theta^G$ is irreducible by Corollary~\ref{cor:15.3}. Characters induced from different $H$-orbits are distinct, since their restrictions to $N$ are sums of disjoint sets of irreducible characters.\\
    Next let $\lambda$ be an irreducible character of $G/N$, and let $\widetilde{\lambda}$ be its lift to $G$. By Lemma~\ref{lem:15.6}, $\widetilde{\lambda}$ is irreducible. Also,
    \[
        \widetilde{\lambda}|_N=\lambda(1)1_N.
    \]
    Thus, for every non-principal $\theta\in\operatorname{Irr}(N)$,
    \[
        (\theta^G,\widetilde{\lambda})_G
        =
        (\theta,\widetilde{\lambda}|_N)_N
        =
        \lambda(1)(\theta,1_N)_N
        =
        0.
    \]
    Hence no character induced from a non-principal character of $N$ is equal to a lifted character.\\
    Choose a set $R$ of representatives for the $H$-orbits on the non-principal irreducible characters of $N$. Since every such orbit has order $|H|$, and conjugation preserves character degrees,
    \[
        |N|-1
        =
        \sum_{\substack{\theta\in\operatorname{Irr}(N)\\\theta\neq1_N}}\theta(1)^2
        =
        |H|\sum_{\theta\in R}\theta(1)^2.
    \]
    The lifted characters have total degree-square sum
    \[
        \sum_{\lambda\in\operatorname{Irr}(G/N)}\widetilde{\lambda}(1)^2
        =
        |G/N|
        =
        |H|,
    \]
    while $\theta^G(1)=|H|\theta(1)$. Therefore the total degree-square sum of the lifted characters and the characters induced from $R$ is
    \begin{align*}
        |H|+\sum_{\theta\in R}\bigl(|H|\theta(1)\bigr)^2
        &=
        |H|+|H|^2\frac{|N|-1}{|H|} \\
        &=
        |H||N|
        =
        |G|.
    \end{align*}
    Corollary~\ref{cor:14.2} now shows that these are all the irreducible characters of $G$. Thus every irreducible character of $G$ is either a lift from $G/N$ or is induced from a non-principal irreducible character of $N$.\\
    Finally, if $\theta,\psi\in\operatorname{Irr}(N)-\{1_N\}$, then
    \[
        (\theta^G,\psi^G)_G
        =
        (\theta,\psi^G|_N)_N
        =
        \sum_{h\in H}(\theta,\psi^h)_N.
    \]
    Since $H$ acts freely, this inner product is $1$ exactly when $\theta=\psi^h$ for some $h\in H$, and it is $0$ otherwise. Both induced characters are irreducible, so
    \[
        \theta^G=\psi^G
        \quad\Longleftrightarrow\quad
        \theta\text{ and }\psi\text{ are $G$-conjugate}.
    \]
    Equivalently, they must lie in the same $H$-orbit. This is also the same as lying in the same $G$-orbit, because $N$ acts trivially on its characters: characters are constant on $N$-conjugacy classes. If $\theta\neq\psi$, the conjugating element of $H$ is necessarily non-identity.
\end{solution}

\subsection*{Further Exercises}

\noindent In this set of exercises, we develop the character table of $G = \mathrm{GL}(2, q)$, where $q$ is any prime power. Recall from Proposition~\ref{prop:4.2} that we have $|G| = q (q-1)^2 (q+1)$. Our first task is to determine the conjugacy classes of $G$. Here we sketch a comparatively ``hands-on'' approach, which uses only some standard topics from linear algebra. (A slicker approach might, for instance, make use of the invariant factor formulation of rational canonical form.) We denote by $m_g(X)$ the minimal polynomial of an element $g \in G$. Our strategy is to exploit the following:

\begin{exercise}
\label{ex:16.8}
Show that conjugate elements of $G$ have the same minimal polynomial.
\end{exercise}
\begin{solution}
    Suppose that $h=ygy^{-1}$ for some $y\in G$. If
    \[
        f(X)=c_0+c_1X+\cdots+c_nX^n\in\mathbb{F}_q[X],
    \]
    then
    \[
        f(h)
        =
        c_0I+c_1ygy^{-1}+\cdots+c_nyg^ny^{-1}
        =
        yf(g)y^{-1}.
    \]
    Hence $f(h)=0$ if and only if $f(g)=0$. Thus $g$ and $h$ are annihilated by exactly the same polynomials. In particular, the monic polynomial of least degree that annihilates $g$ is the same as the monic polynomial of least degree that annihilates $h$. Therefore $m_g(X)=m_h(X)$.
\end{solution}

\noindent Since $m_g(X)$ is a factor of the characteristic polynomial of $g$, which is the quadratic polynomial $\det(XI - g)$, we see that $m_g(X)$ is either linear, or a product of two like or unlike linear factors, or an irreducible monic quadratic.

\begin{exercise}
\label{ex:16.9}
(cont.) Show that if $m_g(X) = X - a$ for some $a \in \mathbb{F}_q^\times$, then $g = \begin{pmatrix} a & 0 \\ 0 & a \end{pmatrix} \in Z(G)$.
\end{exercise}
\begin{solution}
    By definition, the minimal polynomial annihilates $g$. Therefore
    \[
        0=m_g(g)=g-aI,
    \]
    and hence
    \[
        g=aI=
        \begin{pmatrix}
            a&0\\
            0&a
        \end{pmatrix}.
    \]
    A scalar matrix commutes with every element of $G$, so $g\in Z(G)$.
\end{solution}

\begin{exercise}
\label{ex:16.10}
(cont.) Show that if $m_g(X) = (X - a)^2$ for some $a \in \mathbb{F}_q^\times$, then $g$ is conjugate with $\begin{pmatrix} a & 1 \\ 0 & a \end{pmatrix}$, and the conjugacy class of $g$ has order $q^2 - 1$.
\end{exercise}
\begin{solution}
    Put $N=g-aI$. Since $m_g(X)=(X-a)^2$, we have $N^2=0$ and $N\neq0$. Choose $v\in\mathbb{F}_q^2$ such that $Nv\neq0$. The vectors $Nv$ and $v$ are linearly independent: if $v=cNv$, then
    \[
        Nv=cN^2v=0,
    \]
    which is a contradiction. With respect to the ordered basis $(Nv,v)$, the transformation $N$ has matrix
    \[
        \begin{pmatrix}
            0&1\\
            0&0
        \end{pmatrix}.
    \]
    Hence $g=aI+N$ is conjugate to
    \[
        J_a=
        \begin{pmatrix}
            a&1\\
            0&a
        \end{pmatrix}.
    \]
    It remains to compute the centralizer of $J_a$. Since $aI$ is central, a matrix centralizes $J_a$ if and only if it centralizes $J_a-aI$. Direct multiplication shows that
    \[
        \begin{pmatrix}
            u&v\\
            w&z
        \end{pmatrix}
        \begin{pmatrix}
            0&1\\
            0&0
        \end{pmatrix}
        =
        \begin{pmatrix}
            0&1\\
            0&0
        \end{pmatrix}
        \begin{pmatrix}
            u&v\\
            w&z
        \end{pmatrix}
    \]
    if and only if $w=0$ and $z=u$. Thus
    \[
        C_G(J_a)
        =
        \left\{
        \begin{pmatrix}
            u&v\\
            0&u
        \end{pmatrix}
        \,\middle|\,
        u\in\mathbb{F}_q^\times,\ v\in\mathbb{F}_q
        \right\},
    \]
    so $|C_G(J_a)|=q(q-1)$. By Proposition~\ref{prop:3.13}, the conjugacy class of $J_a$ has order
    \[
        |G : C_G(J_a)|
        =
        \frac{q(q-1)^2(q+1)}{q(q-1)}
        =
        q^2-1.
    \]
\end{solution}

\begin{exercise}
\label{ex:16.11}
(cont.) Show that if $m_g(X) = (X - a)(X - b)$ for some distinct $a, b \in \mathbb{F}_q^\times$, then $g$ is conjugate with $\begin{pmatrix} a & 0 \\ 0 & b \end{pmatrix}$, and the conjugacy class of $g$ has order $q^2 + q$.
\end{exercise}
\begin{solution}
    Since $m_g(X)=(X-a)(X-b)$ has two distinct roots, $g$ is diagonalizable over $\mathbb{F}_q$. Choosing eigenvectors for the eigenvalues $a$ and $b$ gives a basis with respect to which
    \[
        g=D_{a,b}=
        \begin{pmatrix}
            a&0\\
            0&b
        \end{pmatrix}.
    \]
    Let
    \[
        M=
        \begin{pmatrix}
            u&v\\
            w&z
        \end{pmatrix}\in G.
    \]
    The equation $MD_{a,b}=D_{a,b}M$ gives
    \[
        (b-a)v=0,
        \qquad
        (a-b)w=0.
    \]
    Since $a\neq b$, we obtain $v=w=0$. Therefore
    \[
        C_G(D_{a,b})
        =
        \left\{
        \begin{pmatrix}
            u&0\\
            0&z
        \end{pmatrix}
        \,\middle|\,
        u,z\in\mathbb{F}_q^\times
        \right\},
    \]
    and so $|C_G(D_{a,b})|=(q-1)^2$. By Proposition~\ref{prop:3.13}, the conjugacy class has order
    \[
        |G : C_G(D_{a,b})|
        =
        \frac{q(q-1)^2(q+1)}{(q-1)^2}
        =
        q^2+q.
    \]
    Interchanging the two eigenvectors shows that $D_{a,b}$ and $D_{b,a}$ are conjugate. Thus these classes are indexed by unordered pairs $\{a,b\}$.
\end{solution}

\begin{exercise}
\label{ex:16.12}
(cont.) Show that if $r, s \in \mathbb{F}_q$ are such that $X^2 - rX - s$ is irreducible, then the conjugacy class of $\begin{pmatrix} 0 & 1 \\ s & r \end{pmatrix}$ has order $q^2 - q$.
\end{exercise}
\begin{solution}
    Let
    \[
        A=
        \begin{pmatrix}
            0&1\\
            s&r
        \end{pmatrix}.
    \]
    Its characteristic polynomial is
    \[
        \det(XI-A)=X^2-rX-s,
    \]
    which is irreducible by hypothesis. Hence the minimal polynomial of $A$ is also $X^2-rX-s$.\\
    Let
    \[
        M=
        \begin{pmatrix}
            u&v\\
            w&z
        \end{pmatrix}.
    \]
    Comparing the entries of $MA$ and $AM$ shows that $MA=AM$ if and only if
    \[
        w=sv,
        \qquad
        z=u+rv.
    \]
    Consequently, the centralizer of $A$ in the algebra of all $2\times2$ matrices over $\mathbb{F}_q$ is
    \[
        \{uI+vA\mid u,v\in\mathbb{F}_q\}
        =
        \mathbb{F}_q[A].
    \]
    Since the minimal polynomial of $A$ is irreducible of degree $2$,
    \[
        \mathbb{F}_q[A]
        \cong
        \mathbb{F}_q[X]/(X^2-rX-s)
    \]
    is a field with $q^2$ elements. Every non-zero element of this field is invertible, and therefore
    \[
        C_G(A)=\mathbb{F}_q[A]^\times
    \]
    has order $q^2-1$. By Proposition~\ref{prop:3.13}, the conjugacy class of $A$ has order
    \[
        |G : C_G(A)|
        =
        \frac{q(q-1)^2(q+1)}{q^2-1}
        =
        q^2-q.
    \]
\end{solution}

\noindent (Hint for Exercises~\ref{ex:16.10}--\ref{ex:16.12}: Compute centralizers and use Proposition~\ref{prop:3.13}.)

\begin{exercise}
\label{ex:16.13}
(cont.) Verify that what follows is a complete list of the conjugacy classes of $G$: $q - 1$ classes of $1$ element each, indexed by elements of $\mathbb{F}_q^\times$; $q - 1$ classes of $q^2 - 1$ elements each, indexed by elements of $\mathbb{F}_q^\times$; $(q - 1)(q - 2)/2$ classes of $q^2 + q$ elements each, indexed by unordered pairs of distinct elements of $\mathbb{F}_q^\times$; and $q(q - 1)/2$ classes of $q^2 - q$ elements each, indexed by irreducible monic quadratic polynomials with coefficients in $\mathbb{F}_q$.
\end{exercise}
\begin{solution}
    Exercises~\ref{ex:16.9}--\ref{ex:16.12} give the four possible types. We check their parameters and then verify that their class sizes account for every element of $G$.\\
    If $m_g(X)=X-a$, Exercise~\ref{ex:16.9} gives the scalar matrix $aI$. Thus there are $q-1$ such classes, each of order $1$.\\
    If $m_g(X)=(X-a)^2$, Exercise~\ref{ex:16.10} gives one class represented by
    \[
        \begin{pmatrix}
            a&1\\
            0&a
        \end{pmatrix}
    \]
    for each $a\in\mathbb{F}_q^\times$. Thus there are $q-1$ such classes, each of order $q^2-1$.\\
    If $m_g(X)=(X-a)(X-b)$ with $a\neq b$, Exercise~\ref{ex:16.11} gives one class for each unordered pair $\{a,b\}\subseteq\mathbb{F}_q^\times$. The number of these pairs is
    \[
        \binom{q-1}{2}
        =
        \frac{(q-1)(q-2)}{2},
    \]
    and each class has order $q^2+q$.\\
    Finally, suppose that $m_g(X)=X^2-rX-s$ is irreducible. In particular, $s\neq0$. Choose $0\neq v\in\mathbb{F}_q^2$. Since $g$ has no eigenvector over $\mathbb{F}_q$, the vectors $v$ and $gv$ are linearly independent. From $g^2=rg+sI$, we see that with respect to the ordered basis $(v,s^{-1}gv)$, the matrix of $g$ is
    \[
        \begin{pmatrix}
            0&1\\
            s&r
        \end{pmatrix}.
    \]
    Thus each irreducible monic quadratic gives exactly one class, and different irreducible polynomials give different classes by Exercise~\ref{ex:16.8}. There are $q^2$ monic quadratic polynomials over $\mathbb{F}_q$. A reducible monic quadratic is determined by an unordered pair of roots in $\mathbb{F}_q$, allowing a repeated root, so there are
    \[
        \binom{q+1}{2}=\frac{q(q+1)}{2}
    \]
    reducible monic quadratics. Hence the number of irreducible monic quadratics is
    \[
        q^2-\frac{q(q+1)}{2}
        =
        \frac{q(q-1)}{2}.
    \]
    By Exercise~\ref{ex:16.12}, each corresponding class has order $q^2-q$.\\
    As a final check, the total number of elements in the listed classes is
    \begin{align*}
        &(q-1)
        +(q-1)(q^2-1)
        +\frac{(q-1)(q-2)}{2}(q^2+q) \\
        &\qquad
        +\frac{q(q-1)}{2}(q^2-q)
        =
        q(q-1)^2(q+1)
        =
        |G|.
    \end{align*}
    Hence no elements or conjugacy classes are missing.
\end{solution}

\noindent Having determined the classes of $G$, we now turn to finding irreducible characters. Recall that we have an epimorphism $\det: G \to \mathbb{F}_q^\times$. Consequently, we can lift characters of $\mathbb{F}_q^\times$ to characters of $G$, as in Lemma~\ref{lem:15.6}. Now $\mathbb{F}_q^\times$ is an abelian group of order $q - 1$, and hence it has $q - 1$ linear characters, which we denote by $\tau_1, \dots, \tau_{q-1}$. (In fact, it is true (see [17, p. 132]) that $\mathbb{F}_q^\times \cong \Z_{q-1}$, and so we can, upon fixing a generator of $\mathbb{F}_q^\times$, completely specify the $\tau_i$.) For each $1 \leq i \leq q - 1$, let $\chi_i = \tau_i \circ \det$. We have
\[
\begin{array}{c|cccc}
 & \begin{pmatrix} a & 0 \\ 0 & a \end{pmatrix} & \begin{pmatrix} a & 1 \\ 0 & a \end{pmatrix} & \begin{pmatrix} a & 0 \\ 0 & b \end{pmatrix} & \begin{pmatrix} 0 & 1 \\ s & r \end{pmatrix} \\
\hline
\chi_i & \tau_i(a)^2 & \tau_i(a)^2 & \tau_i(a) \tau_i(b) & \tau_i(-s)
\end{array}
\]
The characters $\chi_i$ are linear, and $\chi_1$ is indeed the principal character of $G$.

\noindent As in Chapter 2, let $B$, $U$, and $T$ be the subgroups of $G$ consisting, respectively, of all upper triangular, all upper unitriangular, and all diagonal elements of $G$. We have $T \cong \mathbb{F}_q^\times \times \mathbb{F}_q^\times$, and so we see from Example~\ref{ex:15.eg.2} that the irreducible characters of $T$ have the form $\tau_i \times \tau_j$ (after making appropriate identifications). Now $B = U \rtimes T$ by Proposition~\ref{prop:5.1}, so we can lift each character $\tau_i \times \tau_j$ of $T \cong B/U$ to obtain a character $\theta_{ij}$ of $B$. The epimorphism from $B$ to $T$ sends $\begin{pmatrix} x & y \\ 0 & z \end{pmatrix}$ to $\begin{pmatrix} x & 0 \\ 0 & z \end{pmatrix}$, and consequently we have $\theta_{ij}\begin{pmatrix} x & y \\ 0 & z \end{pmatrix} = (\tau_i \times \tau_j)\begin{pmatrix} x & 0 \\ 0 & z \end{pmatrix} = \tau_i(x) \tau_j(z)$. We now wish to consider the induced characters $\theta_{ij}^G$.

\begin{exercise}
\label{ex:16.14}
(cont.) Show that
\[
\begin{array}{c|cccc}
 & \begin{pmatrix} a & 0 \\ 0 & a \end{pmatrix} & \begin{pmatrix} a & 1 \\ 0 & a \end{pmatrix} & \begin{pmatrix} a & 0 \\ 0 & b \end{pmatrix} & \begin{pmatrix} 0 & 1 \\ s & r \end{pmatrix} \\
\hline
\theta_{ij}^G & (q + 1) \tau_i(a) \tau_j(a) & \tau_i(a) \tau_j(a) & \tau_i(a) \tau_j(b) + \tau_i(b) \tau_j(a) & 0
\end{array}
\]
Observe that $\theta_{ij}^G = \theta_{ji}^G$.
\end{exercise}
\begin{solution}
    We use Proposition~\ref{prop:16.5}. First observe that
    \[
        |B|=q(q-1)^2.
    \]
    For a scalar matrix $aI$, we have $C_G(aI)=G$ and $C_B(aI)=B$. Hence its contribution is multiplied by
    \[
        \frac{|C_G(aI)|}{|C_B(aI)|}=|G : B|=q+1,
    \]
    and therefore
    \[
        \theta_{ij}^G(aI)
        =
        (q+1)\theta_{ij}(aI)
        =
        (q+1)\tau_i(a)\tau_j(a).
    \]
    For
    \[
        J_a=
        \begin{pmatrix}
            a&1\\
            0&a
        \end{pmatrix},
    \]
    Exercise~\ref{ex:16.10} shows that $C_G(J_a)\leq B$, so $C_B(J_a)=C_G(J_a)$. Moreover, every element of the $G$-class of $J_a$ which lies in $B$ has the form
    \[
        \begin{pmatrix}
            a&y\\
            0&a
        \end{pmatrix}
        \qquad(y\neq0),
    \]
    and conjugation by diagonal matrices shows that all these elements are conjugate in $B$. Thus there is one relevant $B$-class, its coefficient in Proposition~\ref{prop:16.5} is $1$, and
    \[
        \theta_{ij}^G(J_a)
        =
        \theta_{ij}(J_a)
        =
        \tau_i(a)\tau_j(a).
    \]
    Now let $a\neq b$. The $G$-class of
    \[
        D_{a,b}=
        \begin{pmatrix}
            a&0\\
            0&b
        \end{pmatrix}
    \]
    meets $B$ in two $B$-classes, represented by $D_{a,b}$ and $D_{b,a}$. Indeed, every upper triangular matrix with distinct diagonal entries $a,b$ is conjugate in $B$ to $D_{a,b}$. Also,
    \[
        C_B(D_{a,b})=C_G(D_{a,b})=T,
    \]
    and similarly for $D_{b,a}$. Both coefficients in Proposition~\ref{prop:16.5} are therefore $1$, giving
    \[
        \theta_{ij}^G(D_{a,b})
        =
        \tau_i(a)\tau_j(b)+\tau_i(b)\tau_j(a).
    \]
    Finally, an upper triangular matrix has a characteristic polynomial that splits over $\mathbb{F}_q$. Hence no element whose minimal polynomial is an irreducible quadratic has a conjugate in $B$, so the induced character has value $0$ on the fourth-column classes. Combining these calculations gives the required table. The formula is symmetric in $i$ and $j$, and hence
    \[
        \theta_{ij}^G=\theta_{ji}^G.
    \]
\end{solution}

\begin{exercise}
\label{ex:16.15}
(cont.) Show that $(\theta_{ij}^G, \theta_{ij}^G) = 1 + \delta_{ij}$ and that $(\theta_{ii}^G, \chi_i) = 1$. Conclude that $\theta_{ii}^G - \chi_i$ is an irreducible character of $G$ for each $i$ and that $\theta_{ij}^G$ is an irreducible character of $G$ whenever $i$ and $j$ are distinct.
\end{exercise}
\begin{solution}
    Put $A=\mathbb{F}_q^\times$ and $n=|A|=q-1$, and write $\Theta_{ij}=\theta_{ij}^G$. We compute the inner product using the class sizes from Exercise~\ref{ex:16.13}. Since all values of the linear characters $\tau_i$ have absolute value $1$, the scalar and non-scalar repeated-eigenvalue classes contribute
    \[
        \sum_{a\in A}\left((q+1)^2+(q^2-1)\right)
        =
        2q(q+1)n.
    \]
    There is no contribution from the irreducible-quadratic classes. For the split classes, let
    \[
        S_{ij}
        =
        \sum_{\substack{\{a,b\}\subseteq A\\a\neq b}}
        \left|
        \tau_i(a)\tau_j(b)+\tau_i(b)\tau_j(a)
        \right|^2.
    \]
    Define the linear character $\rho=\tau_i\overline{\tau_j}$ of $A$. Counting ordered pairs and then dividing by $2$, we obtain
    \[
        2S_{ij}
        =
        \sum_{\substack{a,b\in A\\a\neq b}}
        \left(
        2+\rho(a)\overline{\rho(b)}
        +\overline{\rho(a)}\rho(b)
        \right).
    \]
    By character orthogonality on the abelian group $A$,
    \[
        \sum_{a\in A}\rho(a)
        =
        \begin{cases}
            n,&i=j,\\
            0,&i\neq j.
        \end{cases}
    \]
    Consequently,
    \[
        \sum_{a\neq b}\rho(a)\overline{\rho(b)}
        =
        \left|\sum_{a\in A}\rho(a)\right|^2-n,
    \]
    and the same formula holds for its complex conjugate. It follows that
    \[
        S_{ij}=n(n-2)+n^2\delta_{ij}.
    \]
    Since every split class has order $q^2+q=q(q+1)$, we now have
    \begin{align*}
        |G|(\Theta_{ij},\Theta_{ij})_G
        &=
        2q(q+1)n
        +q(q+1)\left(n(n-2)+n^2\delta_{ij}\right) \\
        &=
        q(q+1)n^2(1+\delta_{ij}) \\
        &=
        |G|(1+\delta_{ij}).
    \end{align*}
    Therefore
    \[
        (\theta_{ij}^G,\theta_{ij}^G)_G=1+\delta_{ij}.
    \]
    Next, if
    \[
        b=
        \begin{pmatrix}
            x&y\\
            0&z
        \end{pmatrix}\in B,
    \]
    then
    \[
        \chi_i(b)
        =
        \tau_i(\det b)
        =
        \tau_i(xz)
        =
        \tau_i(x)\tau_i(z)
        =
        \theta_{ii}(b).
    \]
    Thus $\chi_i|_B=\theta_{ii}$, and Frobenius reciprocity gives
    \[
        (\theta_{ii}^G,\chi_i)_G
        =
        (\theta_{ii},\chi_i|_B)_B
        =
        (\theta_{ii},\theta_{ii})_B
        =
        1.
    \]
    If $i\neq j$, then $\theta_{ij}^G$ is a character of norm $1$, so it is irreducible by Corollary~\ref{cor:15.3}. If $i=j$, the preceding inner product shows that $\chi_i$ occurs once in $\theta_{ii}^G$. Hence
    \[
        \eta_i=\theta_{ii}^G-\chi_i
    \]
    is an actual character, and
    \[
        (\eta_i,\eta_i)_G
        =
        2-1-1+1
        =
        1.
    \]
    Therefore $\theta_{ii}^G-\chi_i$ is irreducible. Repeating the same class calculation with two possibly different pairs gives
    \[
        (\theta_{ij}^G,\theta_{k\ell}^G)_G
        =
        \delta_{ik}\delta_{j\ell}+\delta_{i\ell}\delta_{jk}.
    \]
    Thus the characters $\theta_{ij}^G$ are distinct precisely up to interchanging $i$ and $j$. Also, Frobenius reciprocity gives
    \[
        (\theta_{ii}^G,\chi_j)_G
        =
        (\theta_{ii},\chi_j|_B)_B
        =
        (\theta_{ii},\theta_{jj})_B
        =
        \delta_{ij}.
    \]
    Since $(\chi_i,\chi_j)_G=\delta_{ij}$, we obtain
    \begin{align*}
        (\theta_{ii}^G-\chi_i,\theta_{jj}^G-\chi_j)_G
        &=
        2\delta_{ij}-\delta_{ij}-\delta_{ij}+\delta_{ij} \\
        &=
        \delta_{ij}.
    \end{align*}
    Hence the characters $\theta_{ii}^G-\chi_i$ are pairwise distinct.
\end{solution}

\noindent At this point, we have constructed $q - 1$ distinct irreducible characters of degree $1$ (the $\chi_i$), $q - 1$ distinct irreducible characters of degree $q$ (the $\theta_{ii}^G - \chi_i$), and $(q - 1)(q - 2)/2$ distinct irreducible characters of degree $q + 1$ (the $\theta_{ij}^G$ for $i < j$). Our partial character table of $G$ is now
\[
\begin{array}{c|cccc}
 & \begin{pmatrix} a & 0 \\ 0 & a \end{pmatrix} & \begin{pmatrix} a & 1 \\ 0 & a \end{pmatrix} & \begin{pmatrix} a & 0 \\ 0 & b \end{pmatrix} & \begin{pmatrix} 0 & 1 \\ s & r \end{pmatrix} \\
\hline
\chi_i & \tau_i(a)^2 & \tau_i(a)^2 & \tau_i(a) \tau_i(b) & \tau_i(-s) \\
\theta_{ii}^G - \chi_i & q \tau_i(a)^2 & 0 & \tau_i(a) \tau_i(b) & -\tau_i(-s) \\
\theta_{ij}^G & (q + 1) \tau_i(a) \tau_j(a) & \tau_i(a) \tau_j(a) & \tau_i(a) \tau_j(b) + \tau_i(b) \tau_j(a) & 0
\end{array}
\]

\noindent To complete the character table, it will be convenient to work with a different set of representatives for the conjugacy classes of those elements whose minimal polynomials are irreducible quadratics. (We shall refer to these classes as being the ``fourth-column'' classes, owing to their placement in our partial character table above.) By Exercise~\ref{ex:4.4}, $G$ has a subgroup $C$ which is isomorphic with $\mathbb{F}_{q^2}^\times$, where $\mathbb{F}_{q^2}$ is the field of $q^2$ elements. As we remarked before, we have $C \cong \Z_{q^2 - 1}$. It is true, although we are not in a position to prove it, that any subgroup of $G$ that is isomorphic with $\mathbb{F}_{q^2}^\times$ is a conjugate of $C$. (The proof of this fact requires some familiarity with field theory, but the essential point is that all fields of order $q^2$ are isomorphic; see [19, Corollary XIV.2.7].) It is also true that $C$ must contain $Z(G) \cong \mathbb{F}_q^\times$, and so we will consider $\mathbb{F}_q^\times$ as being a subgroup of $C$.

\begin{exercise}
\label{ex:16.16}
(cont.) Show that the elements of $C - \mathbb{F}_q^\times$ form a double set of representatives of the fourth-column classes, with $x \in C - \mathbb{F}_q^\times$ being conjugate with (but unequal to) $x^q$. (\textsc{Hint}: Argue that each element of $G$ whose minimal polynomial is an irreducible quadratic determines a subgroup of $G$ that is isomorphic with $\mathbb{F}_{q^2}^\times$.)
\end{exercise}
\begin{solution}
    We use the embedding from Exercise~\ref{ex:4.4}, in which $\mathbb{F}_{q^2}$ acts on itself by multiplication as a $2$-dimensional vector space over $\mathbb{F}_q$. Thus
    \[
        C=\{m_x\mid x\in\mathbb{F}_{q^2}^\times\},
    \]
    where $m_x(y)=xy$, and the subgroup $\mathbb{F}_q^\times$ consists exactly of the scalar multiplication maps. We identify $m_x$ with $x$.\\
    Let $x\in C-\mathbb{F}_q^\times$. Since $x\notin\mathbb{F}_q$, its minimal polynomial over $\mathbb{F}_q$ is irreducible of degree $2$. The Frobenius map
    \[
        \sigma\colon\mathbb{F}_{q^2}\to\mathbb{F}_{q^2},
        \qquad
        y\mapsto y^q,
    \]
    is an invertible $\mathbb{F}_q$-linear map, and hence represents an element of $G$. For every $y\in\mathbb{F}_{q^2}$,
    \[
        \sigma m_x\sigma^{-1}(y)=x^qy.
    \]
    Therefore $\sigma x\sigma^{-1}=x^q$, so $x$ and $x^q$ are conjugate in $G$. They are unequal, since the elements fixed by Frobenius are exactly the elements of $\mathbb{F}_q$.\\
    Conversely, suppose that $x,y\in C-\mathbb{F}_q^\times$ are conjugate in $G$. By Exercise~\ref{ex:16.8}, they have the same minimal polynomial. The two roots of the minimal polynomial of $x$ in $\mathbb{F}_{q^2}$ are $x$ and $x^q$. Hence
    \[
        y=x
        \qquad\text{or}\qquad
        y=x^q.
    \]
    Thus every $G$-class meets $C-\mathbb{F}_q^\times$ in at most the two elements $x,x^q$.\\
    It remains to show that every fourth-column class meets $C$. Let $g\in G$ have irreducible quadratic minimal polynomial. Then
    \[
        \mathbb{F}_q[g]\cong\mathbb{F}_{q^2},
    \]
    so $\mathbb{F}_q[g]^\times\leq G$ is a subgroup isomorphic to $\mathbb{F}_{q^2}^\times$ and containing $g$. By the conjugacy fact about such subgroups stated immediately before this exercise, $\mathbb{F}_q[g]^\times$ is conjugate to $C$. Hence $g$ is conjugate to an element of $C-\mathbb{F}_q^\times$. Therefore every fourth-column class is represented by exactly two elements of $C-\mathbb{F}_q^\times$, namely $x$ and $x^q$.
\end{solution}

\noindent We now take $x \in C - \mathbb{F}_q^\times$ as our arbitrary fourth-column class representative. We have $\chi_i(x) = \tau_i(\det x)$, $(\theta_{ii}^G - \chi_i)(x) = -\tau_i(\det x)$, and $\theta_{ij}^G(x) = 0$.

\noindent Let $\beta_1, \dots, \beta_{q^2 - 1}$ be the irreducible characters of $C$. (By choosing a generator for $C$, we can specify the $\beta_i$.) We will consider the induced characters $\beta_k^G$.

\begin{exercise}
\label{ex:16.17}
(cont.) Show
\[
\begin{array}{c|cccc}
 & \begin{pmatrix} a & 0 \\ 0 & a \end{pmatrix} & \begin{pmatrix} a & 1 \\ 0 & a \end{pmatrix} & \begin{pmatrix} a & 0 \\ 0 & b \end{pmatrix} & x \in C - \mathbb{F}_q^\times \\
\hline
\beta_k^G & (q^2 - q) \beta_k(a) & 0 & 0 & \beta_k(x) + \beta_k(x)^q
\end{array}
\]
Observe that $(\beta_k^q)^G = \beta_k^G$, where by $\beta_k^q$ we mean the linear character of $C$ defined by $\beta_k^q(x) = \beta_k(x)^q$.
\end{exercise}
\begin{solution}
    We apply Proposition~\ref{prop:16.5} with the subgroup $C$. Since $C$ is abelian, every conjugacy class in $C$ has one element.\\
    If $a\in\mathbb{F}_q^\times$, then $aI\in Z(G)$, so
    \[
        C_G(aI)=G,
        \qquad
        C_C(aI)=C.
    \]
    Hence
    \[
        \beta_k^G(aI)
        =
        \frac{|G|}{|C|}\beta_k(a)
        =
        (q^2-q)\beta_k(a).
    \]
    A non-scalar element of $C$ has irreducible quadratic minimal polynomial. Therefore no element of $C$ is conjugate to a non-scalar repeated-eigenvalue matrix or to a split matrix with two distinct eigenvalues. Proposition~\ref{prop:16.5} consequently gives value $0$ on the second and third types.\\
    Now let $x\in C-\mathbb{F}_q^\times$. By Exercise~\ref{ex:16.16}, the $G$-class of $x$ meets $C$ in exactly the two elements $x$ and $x^q$. Moreover,
    \[
        C_G(x)=C.
    \]
    Indeed, any $\mathbb{F}_q$-linear map $T$ commuting with multiplication by $x$ commutes with every multiplication map $m_y$ for $y\in\mathbb{F}_q[x]=\mathbb{F}_{q^2}$. Hence
    \[
        T(y)=T(m_y(1))=m_y(T(1))=yT(1),
    \]
    so $T$ is itself multiplication by $T(1)\in\mathbb{F}_{q^2}$. Thus both coefficients in Proposition~\ref{prop:16.5} are $1$, and
    \[
        \beta_k^G(x)
        =
        \beta_k(x)+\beta_k(x^q)
        =
        \beta_k(x)+\beta_k(x)^q.
    \]
    This proves the required table.\\
    Finally, define $\beta_k^q$ by
    \[
        \beta_k^q(y)=\beta_k(y)^q=\beta_k(y^q).
    \]
    If $a\in\mathbb{F}_q^\times$, then $a^q=a$, so $\beta_k^q(a)=\beta_k(a)$. If $x\in C-\mathbb{F}_q^\times$, then
    \[
        (\beta_k^q)^G(x)
        =
        \beta_k(x)^q+\beta_k(x)^{q^2}
        =
        \beta_k(x)^q+\beta_k(x),
    \]
    because every value of $\beta_k$ is a $(q^2-1)$th root of unity. Hence
    \[
        (\beta_k^q)^G=\beta_k^G.
    \]
\end{solution}

\noindent It is not hard to show that there are exactly $q - 1$ values of $k$ for which $\beta_k^q = \beta_k$. (This follows, for instance, from the observation made in Example~\ref{ex:15.eg.1} that the linear characters of a cyclic group of order $n$ themselves form a cyclic group of order $n$.)

\noindent Choose some $k$ such that $\beta_k^q \neq \beta_k$, and let $i_k$ be such that the restriction of $\beta_k$ to $\mathbb{F}_q^\times \subseteq C$ is $\tau_{i_k}$. Let $\psi_k = (\theta_{11}^G - \chi_1) \theta_{i_k 1}^G - \theta_{i_k 1}^G - \beta_k^G$; this is a virtual character.

\begin{exercise}
\label{ex:16.18}
(cont.) Show that
\[
\begin{array}{c|cccc}
 & \begin{pmatrix} a & 0 \\ 0 & a \end{pmatrix} & \begin{pmatrix} a & 1 \\ 0 & a \end{pmatrix} & \begin{pmatrix} a & 0 \\ 0 & b \end{pmatrix} & x \in C - \mathbb{F}_q^\times \\
\hline
\psi_k & (q - 1) \tau_{i_k}(a) & -\tau_{i_k}(a) & 0 & -\beta_k(x) - \beta_k(x)^q
\end{array}
\]
\end{exercise}
\begin{solution}
    Put $i=i_k$. From the preceding tables, the three characters occurring in the definition of $\psi_k$ have the following values:
    \[
        \begin{array}{c|cccc}
            &aI&
            \begin{pmatrix}a&1\\0&a\end{pmatrix}&
            \begin{pmatrix}a&0\\0&b\end{pmatrix}&
            x\in C-\mathbb{F}_q^\times\\
            \hline
            \theta_{11}^G-\chi_1
            &q&0&1&-1\\
            \theta_{i1}^G
            &(q+1)\tau_i(a)&\tau_i(a)&\tau_i(a)+\tau_i(b)&0\\
            \beta_k^G
            &(q^2-q)\tau_i(a)&0&0&\beta_k(x)+\beta_k(x)^q
        \end{array}.
    \]
    Here we used $\beta_k|_{\mathbb{F}_q^\times}=\tau_i$. Since
    \[
        \psi_k
        =
        (\theta_{11}^G-\chi_1)\theta_{i1}^G
        -\theta_{i1}^G
        -\beta_k^G,
    \]
    its value at a scalar matrix is
    \begin{align*}
        \psi_k(aI)
        &=
        q(q+1)\tau_i(a)
        -(q+1)\tau_i(a)
        -(q^2-q)\tau_i(a) \\
        &=
        (q-1)\tau_i(a).
    \end{align*}
    At a non-scalar repeated-eigenvalue matrix, we obtain
    \[
        \psi_k
        \begin{pmatrix}
            a&1\\
            0&a
        \end{pmatrix}
        =
        0-\tau_i(a)-0
        =
        -\tau_i(a).
    \]
    At a split matrix with $a\neq b$,
    \[
        \psi_k
        \begin{pmatrix}
            a&0\\
            0&b
        \end{pmatrix}
        =
        \bigl(\tau_i(a)+\tau_i(b)\bigr)
        -\bigl(\tau_i(a)+\tau_i(b)\bigr)
        =
        0.
    \]
    Finally, if $x\in C-\mathbb{F}_q^\times$, then
    \[
        \psi_k(x)=-\beta_k(x)-\beta_k(x)^q.
    \]
    These are precisely the required values.
\end{solution}

\begin{exercise}
\label{ex:16.19}
(cont.) If $k$ is such that $\beta_k^q \neq \beta_k$, show that $(\psi_k, \psi_k) = 1$. Conclude that we have constructed $q(q - 1)/2$ distinct irreducible characters of degree $q - 1$.
\end{exercise}
\begin{solution}
    Write $\beta=\beta_k$. The hypothesis $\beta^q\neq\beta$ says that the character
    \[
        \gamma=\beta^{q-1}
    \]
    is non-principal on $C$. Hence
    \[
        \sum_{x\in C}\gamma(x)=0.
    \]
    If $a\in\mathbb{F}_q^\times$, then $a^{q-1}=1$, and therefore $\gamma(a)=1$. It follows that
    \[
        \sum_{x\in C-\mathbb{F}_q^\times}\beta(x)^{q-1}
        =
        -(q-1).
    \]
    The inverse character $\gamma^{-1}=\beta^{1-q}$ is also non-principal, so similarly
    \[
        \sum_{x\in C-\mathbb{F}_q^\times}\beta(x)^{1-q}
        =
        -(q-1).
    \]
    Since character values have absolute value $1$, we have
    \begin{align*}
        \left|\beta(x)+\beta(x)^q\right|^2
        &=
        \left(\beta(x)+\beta(x)^q\right)
        \left(\overline{\beta(x)}+\overline{\beta(x)^q}\right) \\
        &=
        2+\beta(x)^{q-1}+\beta(x)^{1-q}.
    \end{align*}
    Therefore
    \[
        \sum_{x\in C-\mathbb{F}_q^\times}
        \left|\beta(x)+\beta(x)^q\right|^2
        =
        2(q^2-q)-2(q-1)
        =
        2(q-1)^2.
    \]
    Each fourth-column class has order $q^2-q$ and is represented twice in $C-\mathbb{F}_q^\times$. Using the values from Exercise~\ref{ex:16.18}, we obtain
    \begin{align*}
        |G|(\psi_k,\psi_k)_G
        &=
        (q-1)^3
        +(q^2-1)(q-1) \\
        &\qquad
        +\frac{q^2-q}{2}
        \sum_{x\in C-\mathbb{F}_q^\times}
        \left|\beta(x)+\beta(x)^q\right|^2 \\
        &=
        (q-1)^3
        +(q^2-1)(q-1)
        +(q^2-q)(q-1)^2 \\
        &=
        q(q-1)^2(q+1)
        =
        |G|.
    \end{align*}
    Thus $(\psi_k,\psi_k)_G=1$. Since $\psi_k$ is a virtual character, it is an integral linear combination of the irreducible characters of $G$. Its norm being $1$ means that it is either an irreducible character or the negative of one. However,
    \[
        \psi_k(1)=q-1>0,
    \]
    so $\psi_k$ is an irreducible character of degree $q-1$.\\
    The character group of the cyclic group $C$ is cyclic of order $q^2-1$. A character $\beta$ satisfies $\beta^q=\beta$ exactly when $\beta^{q-1}=1$, and there are $q-1$ such characters. Thus there are $q^2-q$ characters for which $\beta^q\neq\beta$. The operation $\beta\mapsto\beta^q$ is an involution, since $\beta^{q^2}=\beta$, so these characters form
    \[
        \frac{q^2-q}{2}
    \]
    pairs $\{\beta,\beta^q\}$. Exercise~\ref{ex:16.17} shows that $\beta^G=(\beta^q)^G$, and the restrictions of $\beta$ and $\beta^q$ to $\mathbb{F}_q^\times$ are equal. Hence the two characters in each pair give the same $\psi_k$.\\
    Conversely, write $\psi_\beta$ for the character constructed from $\beta$, and suppose that $\psi_\beta=\psi_\delta$. Comparing scalar values shows that $\beta$ and $\delta$ have the same restriction to $\mathbb{F}_q^\times$. Comparing the fourth-column values then gives
    \[
        \beta+\beta^q=\delta+\delta^q
    \]
    on all of $C$. The irreducible characters of the abelian group $C$ are linearly independent, so $\{\delta,\delta^q\}=\{\beta,\beta^q\}$. Therefore distinct Frobenius pairs give distinct irreducible characters. The number constructed is consequently
    \[
        \frac{q^2-q}{2}
        =
        \frac{q(q-1)}{2}.
    \]
\end{solution}

\noindent The number of irreducible characters that we have constructed is now equal to the number of conjugacy classes of $G$, and so by Theorem~\ref{thm:14.3} we are done.

\end{document}
